\section{泛函的表示}
\label{sec:lfunc-representation}

上一节逐个研究读数规则。现在把同一空间上全部连续线性读数收集起来，看它们组成怎样的空间，以及在具体空间中它们能否写成熟悉的形式。先建立对偶空间及其范数，并证明它总是完备的；再分别在Hilbert空间和$L^p$空间中证明，每个连续线性泛函都是内积或加权积分，表示唯一，并且泛函范数等于表示对象的范数。

\subsection{对偶空间}
\label{sec:lfunc-dual-space}

两个有界线性泛函之和、数乘仍是有界线性泛函，所以它们组成一个线性空间，泛函范数是这个空间上的范数：三角不等式来自$\lvert(\ell_1+\ell_2)(x)\rvert\leq\lvert\ell_1(x)\rvert+\lvert\ell_2(x)\rvert$再对单位球取上确界。

\begin{definition}[对偶空间与二次对偶空间]
\label{def:lfunc-dual-space}
赋范空间$X$上全部有界线性泛函组成的空间记为$X^*$，配以泛函范数，称为$X$的\term{对偶空间}（dual space）。$X^*$本身是赋范空间，它的对偶空间$X^{**}=(X^*)^*$称为$X$的\term{二次对偶空间}（bidual space）。
\end{definition}

对偶空间的元素是读数规则，二次对偶空间的元素是“读数规则的读数规则”：给它一个$X$上的泛函，它返回一个数。第~\ref{sec:lfunc-reflexivity}小节会看到，每个$x\in X$都以$\ell\mapsto\ell(x)$的方式自然地成为$X^{**}$的元素；比较这些元素与$X^{**}$的全部元素，是后面讨论弱拓扑与紧性的基础。

\paragraph{对偶范数可以不同于原范数}
即使在有限维中，对偶空间的范数也可以与原空间的范数形式不同。在$\mathbb R^2$上配$\lVert\bm x\rVert_1=\lvert x_1\rvert+\lvert x_2\rvert$。每个线性泛函都写成$\ell(\bm x)=a_1x_1+a_2x_2$，由
\[
\lvert\ell(\bm x)\rvert\leq\max(\lvert a_1\rvert,\lvert a_2\rvert)\bigl(\lvert x_1\rvert+\lvert x_2\rvert\bigr)
\]
且在$\bm x=\bm e_1$或$\bm e_2$处取到等号，得$\lVert\ell\rVert=\max(\lvert a_1\rvert,\lvert a_2\rvert)$。系数向量$\bm a$作为读数的长度要用最大值范数计算，这是\term{对偶范数}（dual norm）的一个例子。所以不能把读数规则与表示它的系数向量随意等同，除非已经核对了范数。

\paragraph{对偶空间总是完备的}
对偶空间有一个原空间未必具有的性质：无论$X$本身是否完备，$X^*$都完备。原因是读数取值在完备的$\mathbb F$中。

\begin{theorem}[对偶空间的完备性]
\label{thm:lfunc-dual-complete}
对任意赋范空间$X$，$X^*$是Banach空间。
\end{theorem}

\begin{proof}
设$(\ell_n)$是$X^*$中的Cauchy列。对每个$x$，$\lvert\ell_n(x)-\ell_m(x)\rvert\leq\lVert\ell_n-\ell_m\rVert\lVert x\rVert$，所以$(\ell_n(x))$是$\mathbb F$中的Cauchy列，有极限，记为$\ell(x)$。极限保持线性，$\ell$是线性泛函。

给定$\varepsilon>0$，取$N$使$n,m\geq N$时$\lVert\ell_n-\ell_m\rVert<\varepsilon$。对$\lVert x\rVert\leq1$，$\lvert\ell_n(x)-\ell_m(x)\rvert<\varepsilon$；令$m\to\infty$得$\lvert\ell_n(x)-\ell(x)\rvert\leq\varepsilon$。所以$\lvert\ell(x)\rvert\leq\lVert\ell_N\rVert+\varepsilon$在单位球上成立，$\ell$有界，属于$X^*$；并且$n\geq N$时$\lVert\ell_n-\ell\rVert\leq\varepsilon$，即$\ell_n\to\ell$。
\end{proof}

例如由多项式组成的不完备空间配上一致范数，它的对偶空间仍是Banach空间。这一性质在第~\ref{sec:lfunc-uniform-boundedness}节和第~\ref{sec:lfunc-weak-norm-control}小节中起关键作用：需要完备性的论证，可以在$X^*$上进行而不必要求$X$完备。

\subsection{Hilbert空间的泛函表示}
\label{sec:lfunc-riesz}

在Hilbert空间$H$中，每个向量都给出一个读数。按第~\ref{chap:inner-product-spaces}章的约定，复内积对第一个变量共轭线性、对第二个变量线性。对$h\in H$令$\ell_h(f)=\langle h,f\rangle$，它对$f$线性；由Cauchy--Schwarz不等式，$\lvert\ell_h(f)\rvert\leq\lVert h\rVert\lVert f\rVert$，所以$\lVert\ell_h\rVert\leq\lVert h\rVert$。$h\neq0$时取$f=h/\lVert h\rVert$，读数为$\lVert h\rVert$，于是$\lVert\ell_h\rVert=\lVert h\rVert$。

反过来的问题更有用：给定连续线性泛函$\ell$，是否一定有向量$h$使$\ell=\ell_h$？若是，读数就可以用空间中的几何来研究，读数的范数就是一个向量的长度。

\paragraph{零空间与正交方向}
零泛函由零向量表示。设$\ell\neq0$，考察它的零空间$N=\{f:\ell(f)=0\}$。它是线性子空间；又由连续性，若$f_n\in N$且$f_n\to f$，则$\ell(f)=\lim\ell(f_n)=0$，所以$N$是闭子空间，且$N\neq H$。若能找到与$N$正交的非零向量$w$，读数就只取决于$f$沿$w$方向的分量，表示向量应当是$w$的某个倍数。由此得到的结论称为\term{Riesz表示定理}（Riesz representation theorem）。

\begin{theorem}[Riesz表示定理]
\label{thm:lfunc-riesz}
设$H$是Hilbert空间，$\ell\in H^*$。则存在唯一的$h_\ell\in H$，使
\begin{equation}
\ell(f)=\langle h_\ell,f\rangle\qquad(f\in H),
\label{eq:lfunc-riesz}
\end{equation}
并且$\lVert\ell\rVert=\lVert h_\ell\rVert$。
\end{theorem}

\begin{proof}
$\ell=0$时取$h_\ell=0$。下设$\ell\neq0$，$N$如上。

\emph{找与零空间正交的方向。}取$v$使$\ell(v)\neq0$。由第~\ref{chap:inner-product-spaces}章的正交投影定理，$v=Pv+w$，其中$Pv\in N$，$w\perp N$。由$\ell(w)=\ell(v)-\ell(Pv)=\ell(v)\neq0$知$w\neq0$。

\emph{写出表示。}对任意$f\in H$，向量$f-\frac{\ell(f)}{\ell(w)}w$在$\ell$下的值为零，属于$N$，因而与$w$正交：
\[
0=\Bigl\langle w,f-\frac{\ell(f)}{\ell(w)}w\Bigr\rangle=\langle w,f\rangle-\frac{\ell(f)}{\ell(w)}\lVert w\rVert^2 .
\]
这里用到内积对第二个变量线性。解出$\ell(f)=\frac{\ell(w)}{\lVert w\rVert^2}\langle w,f\rangle$。令$h_\ell=\frac{\overline{\ell(w)}}{\lVert w\rVert^2}w$，由内积对第一个变量共轭线性，$\langle h_\ell,f\rangle=\frac{\ell(w)}{\lVert w\rVert^2}\langle w,f\rangle=\ell(f)$。

\emph{唯一性与范数。}若$h,h'$都满足式~\eqref{eq:lfunc-riesz}，则$\langle h-h',f\rangle=0$对一切$f$成立，取$f=h-h'$得$h=h'$。范数等式是本小节开头对$\ell_h$的计算。
\end{proof}

证明只用了正交投影定理，而正交投影定理对任意Hilbert空间中的闭子空间成立，所以Riesz表示定理不要求$H$可分，也不需要事先选取正交基。表示向量公式中的共轭使对应$\ell\mapsto h_\ell$在复空间中是共轭线性的：$h_{\alpha\ell}=\bar\alpha h_\ell$。实空间中这个共轭消失，$H^*$与$H$通过保范数的线性双射等同。

\paragraph{\texorpdfstring{$L^2$}{L2}中的读数与方向}
在$L^2(0,1)$中，几个读数的表示向量可以直接写出。整体平均$\ell(f)=\int_0^1f$对应$h_\ell=1$，范数为$1$。加权读数$\ell(f)=\int_0^1tf(t)\,dt$对应$h_\ell(t)=t$，范数为$(\int_0^1t^2\,dt)^{1/2}=1/\sqrt3$。第$k$个正交系数$f\mapsto\langle e_k,f\rangle$对应$h_\ell=e_k$。

\paragraph{表示以连续性为前提}
定理不能替一条读数规则证明连续性。点评估在$L^2(0,1)$上根本没有定义；在$(C([0,1]),\lVert\cdot\rVert_2)$上它有定义但不连续，也确实没有表示向量：若有$h$使$f(0)=\langle h,f\rangle$对一切连续$f$成立，代入尖峰$\tau_n$，左端恒为$1$，右端不超过$\lVert h\rVert_2\lVert\tau_n\rVert_2\to0$，矛盾。完备性也不能省：$(C([0,1]),\lVert\cdot\rVert_2)$上的有界读数$\int_0^{1/2}f$在完备化$L^2(0,1)$中的表示向量是$\mathbf{1}_{[0,1/2]}$，它没有连续代表，所以在不完备的空间中找不到表示向量。

若某个函数空间配上内积后完备，并且点评估在其中连续，Riesz表示定理就把点评估写成与某个固定函数的内积。第~\ref{chap:rkhs}章研究的再生核Hilbert空间正是这样的空间。

\subsection{\texorpdfstring{$L^p$}{Lp}空间的泛函表示}
\label{sec:lfunc-lp-dual}

在$L^p$中没有内积可用，自然的候选是加权积分~\eqref{eq:lfunc-weighted-integral}。第~\ref{sec:lfunc-examples}小节已说明每个$g\in L^q$给出有界读数$L_g$，且$\lVert L_g\rVert\leq\lVert g\rVert_q$。本小节回答反方向的问题。为了不处理共轭记号，下面在实值函数上证明；复值情形对实部与虚部分别处理，结论相同，只是$p=2$时$L_g$的表示向量在内积记号下是$\bar g$。

\begin{theorem}[$L^p$的对偶表示]
\label{thm:lfunc-lp-dual}
设$(\Omega,\mathcal F,\mu)$是$\sigma$有限测度空间，$1\leq p<\infty$，$q$为共轭指数。则每个$\ell\in(L^p(\mu))^*$都存在$g\in L^q(\mu)$使$\ell=L_g$；$g$在几乎处处相等的意义下唯一，并且$\lVert\ell\rVert=\lVert g\rVert_q$。
\end{theorem}

证明分三步：先说明$\lVert L_g\rVert$恰好等于$\lVert g\rVert_q$，从而表示唯一；再在有限测度下由读数构造测度并用Radon--Nikodym定理得到$g$；最后推广到$\sigma$有限测度。

\paragraph{范数等式与唯一性}
设$1<p<\infty$，$g\neq0$。令$f=\lvert g\rvert^{q-1}\operatorname{sgn}g$，则$fg=\lvert g\rvert^q$，又因$(q-1)p=q$，$\lvert f\rvert^p=\lvert g\rvert^q$。于是$L_g(f)=\lVert g\rVert_q^q$，$\lVert f\rVert_p=\lVert g\rVert_q^{q/p}$，比值为$\lVert g\rVert_q^{q-q/p}=\lVert g\rVert_q$。$p=1$、$q=\infty$时等号未必取到，改用逼近：对$0<\varepsilon<\lVert g\rVert_\infty$，集合$E=\{\lvert g\rvert>\lVert g\rVert_\infty-\varepsilon\}$测度为正；由$\sigma$有限，可取$E$中测度有限且为正的部分$E'$，令$f=\operatorname{sgn}(g)\mathbf{1}_{E'}/\mu(E')$，则$\lVert f\rVert_1=1$而$L_g(f)\geq\lVert g\rVert_\infty-\varepsilon$。两种情形都得到$\lVert L_g\rVert=\lVert g\rVert_q$。

由此表示是唯一的：若$L_g=L_{g'}$，则$L_{g-g'}=0$，其范数$\lVert g-g'\rVert_q=0$，即$g=g'$几乎处处成立。表示函数只确定到几乎处处相等，这与$L^q$的元素本来就是等价类一致。

\paragraph{有限测度下构造表示函数}
先设$\mu(\Omega)<\infty$。给定$\ell\in(L^p)^*$，对可测集$E$令$\nu(E)=\ell(\mathbf{1}_E)$。有限可加性来自线性。可数可加性来自连续性，这里用到$p<\infty$：若$E=\bigcup_kE_k$为不交并，则
\[
\Bigl\lVert\mathbf{1}_E-\sum_{k\leq n}\mathbf{1}_{E_k}\Bigr\rVert_p=\mu\Bigl(\bigcup_{k>n}E_k\Bigr)^{1/p}\to0,
\]
所以$\sum_{k\leq n}\nu(E_k)\to\nu(E)$。又$\lvert\nu(E)\rvert\leq\lVert\ell\rVert\mu(E)^{1/p}$，$\nu$是有限的有符号测度，并且$\mu(E)=0$时$\nu(E)=0$，即$\nu$对$\mu$绝对连续。

由第~\ref{chap:measure}章，有限的有符号测度可分解为两个有限测度之差，它们都对$\mu$绝对连续；分别应用第~\ref{chap:measure-densities}章的Radon--Nikodym定理，得到$g\in L^1(\mu)$使$\ell(\mathbf{1}_E)=\int_Eg\,d\mu$对一切可测集$E$成立。由线性，$\ell(f)=\int fg\,d\mu$对一切简单函数$f$成立。

\paragraph{表示函数属于\texorpdfstring{$L^q$}{Lq}}
设$1<p<\infty$。令$E_n=\{\lvert g\rvert\leq n\}$，$f_n=\lvert g\rvert^{q-1}\operatorname{sgn}(g)\mathbf{1}_{E_n}$。$f_n$是有界可测函数，可以被简单函数一致逼近；在有限测度下一致收敛推出$L^p$收敛，而$\ell$连续、$g$可积，所以等式$\ell(f_n)=\int f_ng\,d\mu$也成立。于是
\[
\int_{E_n}\lvert g\rvert^q\,d\mu=\ell(f_n)\leq\lVert\ell\rVert\lVert f_n\rVert_p=\lVert\ell\rVert\Bigl(\int_{E_n}\lvert g\rvert^q\,d\mu\Bigr)^{1/p},
\]
整理得$(\int_{E_n}\lvert g\rvert^q)^{1/q}\leq\lVert\ell\rVert$，由单调收敛定理令$n\to\infty$得$\lVert g\rVert_q\leq\lVert\ell\rVert$。

设$p=1$。若$F=\{g>\lVert\ell\rVert\}$测度为正，取$f=\mathbf{1}_F/\mu(F)$，则$\lVert f\rVert_1=1$，而$\ell(f)=\frac{1}{\mu(F)}\int_Fg\,d\mu>\lVert\ell\rVert$，矛盾；对$\{g<-\lVert\ell\rVert\}$同理。所以$\lVert g\rVert_\infty\leq\lVert\ell\rVert$。

两种情形下$g\in L^q$，$\ell$与$L_g$都是$L^p$上的连续线性泛函，在简单函数上相等。由第~\ref{chap:normed-spaces}章，简单函数在$L^p$（$p<\infty$）中稠密，两个连续函数在稠密集上相等则处处相等，所以$\ell=L_g$。

\paragraph{推广到\texorpdfstring{$\sigma$}{sigma}有限测度}
设$\Omega=\bigcup_k\Omega_k$，$\Omega_k$两两不交且测度有限。把$\ell$限制到在$\Omega_k$外为零的函数上，由有限测度情形得到$\Omega_k$上的$g_k$，令$g=\sum_kg_k\mathbf{1}_{\Omega_k}$。对$A_n=\Omega_1\cup\cdots\cup\Omega_n$，同样的论证给出$\lVert g\mathbf{1}_{A_n}\rVert_q\leq\lVert\ell\rVert$，由单调收敛得$\lVert g\rVert_q\leq\lVert\ell\rVert$。于是$\ell$与$L_g$在支撑于某个$A_n$内的简单函数上相等，这些函数在$L^p$中仍稠密，所以$\ell=L_g$。$p=1$时$\sigma$有限的假设不能去掉；$1<p<\infty$时可以用更细的论证去掉，这里不展开。

\begin{example}[一个可复算的对偶范数]
\label{ex:lfunc-l3-norm}
在$L^3(0,1)$上取$\ell(f)=\int_0^1tf(t)\,dt$。共轭指数$q=3/2$，由定理~\ref{thm:lfunc-lp-dual}，
\[
\lVert\ell\rVert=\Bigl(\int_0^1t^{3/2}\,dt\Bigr)^{2/3}=\Bigl(\frac25\Bigr)^{2/3}\approx0.543,
\]
取到范数的输入是$f(t)=t^{q-1}=t^{1/2}$：$\ell(t^{1/2})=\int_0^1t^{3/2}=2/5$，$\lVert t^{1/2}\rVert_3=(\int_0^1t^{3/2})^{1/3}=(2/5)^{1/3}$，比值为$(2/5)^{2/3}$。同一个权重在$L^2$中的读数范数是$\lVert t\rVert_2=1/\sqrt3\approx0.577$，在$L^1$中是$\lVert t\rVert_\infty=1$。
\end{example}

\paragraph{端点\texorpdfstring{$p=\infty$}{p=infinity}}
定理排除了$p=\infty$，这不是技术上的疏忽。每个$g\in L^1$仍给出$L^\infty$上的有界读数$L_g$，但$L^\infty$上的连续线性泛函一般不能全部写成这种形式。第~\ref{sec:lfunc-hahn-banach}小节将用延拓定理在$\ell^\infty$上构造一个不来自任何$\ell^1$序列的连续线性泛函。所以有限维和$L^2$中“对偶就是自身”的直觉，以及$1\leq p<\infty$时“对偶就是$L^q$”的结论，都不能直接用到$L^\infty$上。
